% Part: exercises
% Chapter: first-order-logic
% Section: semantics-identity

\documentclass[../../../../include/open-logic-section]{subfiles}

\begin{document}

\olfileid{exe}{fol}{sid}

\olsection{Semantics with identity}

\begin{prob}[Expressing cardinalities]
Cardinalities are number of things: one things, a thousand things, infinitely many things, as many things as there are natural numbers, as many things as there are real numbers, etc. First-order logic with identity can express some facts 
about cardinalities. 
\begin{enumerate}
	\item Write !!a{formula} of first-order logic that holds only if the domain contains at least two individuals.
	\item Write !!a{formula} of first-order logic that holds only if the domain contains exactly two individuals.
	\item Write !!a{formula} of first-order logic that holds only if the domain contains at most three individuals.
	\item Write !!a{formula} of first-order logic that holds only if the domain contains infinitely many things.
\end{enumerate}
\end{prob}

\begin{prob}[Formalizing claims about quantity]
Write !!{formula}s of first-order logic with identity that express the following: 
\begin{enumerate}
	\item There are at least three things.
	\item There are exactly three things that are red.\\ 
	(Use $\Atom{\Obj F}{x}$ to express the claim that $x$ is red.)
\end{enumerate}

\end{prob}

\begin{prob}[Leibniz's law]
In first-order logic, Leibniz's law is expressed by the following schema, where $!A$ can stand for any formula in which $y$ doesn't appear free:
$$\lforall x \lforall y (\eq[x][y] \lif (!A \liff \Subst{!A}{y}{x}))$$
(Recall that $\Subst{!A}{y}{x}$ stands for the !!{formula} resulting from substituting $y$ for all free occurrences of $x$ in $!A$.) Answer the following.
\begin{enumerate}
	\item Explain why if (for some assignement $g$ in some model $\Struct{M}$) $\Value{y}{M}[g] = \Value{x}{M}[g]$, then $\Sat{M}{\Atom{\Obj F}{y}}[g]$ iff $\Sat{M}{\Atom{\Obj F}{x}}[g]$. 
	\item (Generalize your answer.) Let $\Atom{!A}$ be an atomic !!{formula} in which $y$ doesn't appear free. Explain why, if $\Value{y}{M}[g] = \Value{x}{M}[g]$, $\Sat{M}{!A}[g]$ iff $\Sat{M}{\Subst{!A}{y}{x}}[g]$.
	\item Explain why, if $\Sat{M}{!A}[g]$ iff $\Sat{M}{\Subst{!A}{y}{x}}[g]$ and $\Sat{M}{!B}[g]$ iff $\Sat{M}{\Subst{!B}{y}{x}}[g]$, then $\Sat{M}{!A \land !B}[g]$ and $\Sat{M}{\Subst{!A \land !B}{y}{x}}[g]$.
	\item Suppose that if $\Struct{M}$) $\Value{y}{M}[g] = \Value{x}{M}[g]$, then $\Sat{M}{!A}[g]$ iff $\Sat{M}{\Subst{!A}{y}{x}}[g]$. Explain why, if $\Sat{M}{!A}[g]$ iff $\Sat{M}{\Subst{!A}{y}{x}}[g]$, then $\lforall \nu \Sat{M}{!A}[g]$ iff $\forall \nu \Sat{M}{\Subst{!A}{y}{x}}[g]$ for any variable $\nu$. (Tip: $\nu$ could be $x$ or $y$ themselves, but that doesn't change anything to the explanation of why this hods.)
	\item Using the claims above, explain why Leibniz's law holds in first-order logic with identity. 
\end{enumerate}
\end{prob}

\end{document}
